% Preamble: common macros and styling for HB100 paper

\documentclass[10pt,a4paper]{article}

% Fonts and encoding
\usepackage[utf8]{inputenc}
\usepackage[T1]{fontenc}
\usepackage{lmodern}

% Margins
\usepackage[margin=1in]{geometry}

% Line spacing
\usepackage{setspace}
\onehalfspacing

% Math
\usepackage{amsmath}
\usepackage{amssymb}
\usepackage{mathtools}

% Figures and tables
\usepackage{graphicx}
\usepackage{booktabs}
\usepackage{array}
\usepackage[skip=10pt,labelfont=bf]{caption}

% Color
\usepackage{xcolor}
\definecolor{darkblue}{RGB}{0,51,102}
\definecolor{medblue}{RGB}{0,102,204}

% Hyperlinks
\usepackage{hyperref}
\hypersetup{
    colorlinks=true,
    linkcolor=darkblue,
    citecolor=darkblue,
    urlcolor=medblue,
    pdftitle={HB100 Phased Array Radar},
    pdfauthor={Author}
}

% Bibliography
\usepackage{natbib}
\bibliographystyle{abbrvnat}

% Section formatting
\usepackage{titlesec}
\titleformat{\section}{\large\bfseries\color{darkblue}}{\thesection}{1em}{}
\titleformat{\subsection}{\normalsize\bfseries}{\thesubsection}{1em}{}

% Code listings (for algorithm pseudocode)
\usepackage{listings}
\lstset{
    basicstyle=\ttfamily\small,
    breaklines=true,
    frame=single,
    rulecolor=\color{lightgray}
}

% Spacing tweaks
\setlength{\parindent}{0.5cm}
\setlength{\parskip}{0.3cm}


\title{HB100 Phased Array Doppler Radar: Custom Analog Signal Conditioning for Budget-Friendly Multi-Channel Beamforming}
\author{Anonymous}
\date{\today}

\begin{document}

\maketitle

\begin{abstract}
This paper presents the first documented phased-array radar system using HB100 Doppler microwave modules (10.525 GHz).
We address the challenge of handling free-running Gunn oscillators and bipolar IF outputs through custom analog signal conditioning
(virtual ground biasing, multi-stage op-amp amplification). Our dual processing pipelines --- MUSIC-based angle-of-arrival estimation
and amplitude-only monopulse --- demonstrate robust target detection and localization despite oscillator drift.
We validate performance through simulation and provide open-source hardware, firmware, and signal processing code
to enable reproducibility and community contributions. This work establishes HB100 viability for low-cost phased array research.
\end{abstract}

\section{Introduction}
\label{sec:intro}

The HB100 is a low-cost (\$2--10) Doppler microwave module operating at 10.525~GHz,
widely used in motion detection and security applications. To date, all documented
applications use single-channel configurations. This paper is the first to explore
multi-channel phased-array configurations with HB100 modules.

\subsection{Motivation}

A 4-element uniform linear array (ULA) with 37--38~mm element spacing and phase-coherent
processing enables:
\begin{itemize}
\item Direction-of-arrival (DOA) estimation via MUSIC
\item Multi-target localization
\item Adaptive nulling (future work)
\end{itemize}

Existing budget radar systems cost \$50k--\$200k. With analog signal conditioning,
HB100-based phased arrays could provide research-grade angle estimation at \$500--\$2000.

\subsection{Technical Challenges}

Free-running Gunn oscillators introduce phase drift. The HB100 IF output (10~\textmu V--2~mV,
AC-coupled) cannot be digitized directly by a unipolar ADC. We solve these with:

\begin{enumerate}
\item Virtual ground biasing (1.65~V reference)
\item Multi-stage op-amp gain (1092$\times$) to utilize ADC dynamic range
\item Forward-backward averaging for phase coherence robustness
\item Dual processing pipelines: MUSIC for angle, monopulse for robustness
\end{enumerate}

\subsection{Contributions}

\begin{enumerate}
\item \textbf{Novel system design:} First documented HB100 phased array with complete
  hardware, firmware, and signal processing stack.
\item \textbf{Algorithm comparison:} Quantify MUSIC vs.\ monopulse trade-offs; demonstrate
  $\sim$2--5° RMS angle error despite oscillator drift.
\item \textbf{Open-source reproducibility:} Provide firmware (ESP32-S3), Python DSP code,
  Hydra-managed simulation, and GitHub artifact for community use.
\end{enumerate}

\subsection{Paper Organization}

Section~\ref{sec:related} reviews phased array literature. Section~\ref{sec:design} details
hardware architecture and analog conditioning. Section~\ref{sec:signal} covers MUSIC, monopulse,
and Kalman filtering. Section~\ref{sec:validation} presents simulation results.
Section~\ref{sec:repro} discusses reproducibility, and Section~\ref{sec:conclusion} concludes.

\section{Related Work}
\label{sec:related}

\subsection{Phased Array Signal Processing}

Phased arrays have been extensively studied since the 1960s. We focus on spatial estimation
and direction-of-arrival (DOA) algorithms applicable to small ULAs.

\subsubsection{MUSIC Algorithm}

Schmidt's MUSIC (MUltiple SIgnal Classification) algorithm~\cite{Schmidt1986} estimates
direction-of-arrival by decomposing the spatial correlation matrix into signal and noise subspaces.
For a narrowband signal impinging on a ULA:

\begin{equation}
\mathbf{x}(k) = \mathbf{a}(\theta) s(k) + \mathbf{n}(k)
\end{equation}

where $\mathbf{a}(\theta)$ is the steering vector, $s(k)$ is the target signal, and $\mathbf{n}(k)$
is noise. The MUSIC pseudospectrum is:

\begin{equation}
P_{\text{MUSIC}}(\theta) = \frac{1}{\|\mathbf{a}(\theta)\|_{\mathbf{E}_n}^2}
\end{equation}

where $\mathbf{E}_n$ is the noise subspace (eigenvectors of the smallest eigenvalues).

Spatial smoothing and forward-backward averaging improve performance under correlated noise~\cite{Krim1996}.

\subsubsection{Monopulse DOA}

Monopulse estimation uses amplitude ratios across subarrays~\cite{Johnson1993}:

\begin{equation}
\text{Bias} = \frac{\sum_{\text{right}} |\mathbf{X}| - \sum_{\text{left}} |\mathbf{X}|}{\sum_{\text{right}} |\mathbf{X}| + \sum_{\text{left}} |\mathbf{X}|}
\end{equation}

Monopulse is robust to oscillator phase noise because it ignores phase; it is typically less
accurate than MUSIC but requires no correlation matrix computation.

\subsection{HB100 Doppler Modules}

The HB100 is a low-cost Gunn oscillator module with an integrated mixer for Doppler detection.
Published work uses single-channel configurations for motion detection~\cite{HB100Datasheet},
vehicle counting, and security. To our knowledge, no prior work explores multi-channel phased arrays.

\subsection{Budget Radar Systems}

Amateur radio and hobby communities have built phased arrays with RTL-SDR dongles~\cite{VanTrees2002},
but these are passive and require strong external transmitters. Active phased arrays with integrated
transmitters remain expensive.

Our work bridges this gap by providing a complete, open-source solution.

\section{System Design}
\label{sec:design}

\subsection{Hardware Architecture}

\subsubsection{Array Geometry}

Our system uses a 4-element uniform linear array (ULA) with element spacing 37--38~mm
(center-to-center):

\begin{itemize}
\item Element 0: 0.000~m (GPIO 7)
\item Element 1: 0.037~m (GPIO 6)
\item Element 2: 0.075~m (GPIO 5)
\item Element 3: 0.112~m (GPIO 4)
\end{itemize}

At 10.525~GHz, the wavelength is $\lambda \approx 28.5$~mm, giving a maximum baseline of
$d_{\max} = 112$~mm $\approx 3.93\lambda$. This baseline supports DOA estimation over $\pm 80°$
without aliasing~\cite{Monzingo2011}.

\subsubsection{Analog Signal Conditioning}

Each HB100 IF output (10~\textmu V--2~mV at $\sim$200~Hz nominal Doppler) feeds a custom
op-amp circuit:

\begin{enumerate}
\item \textbf{Virtual Ground Reference:} Two 10~k$\Omega$ resistors divide the 3.3~V ADC rail
  to establish a 1.65~V midpoint reference. This allows the ADC to digitize bipolar IF swings.

\item \textbf{Pre-Amplifier EMI Filter:} 10~\textmu F capacitor in parallel with 10~k$\Omega$
  resistor to virtual ground. Cutoff frequency $f_c \approx 1.6$~Hz attenuates 50/60~Hz powerline hum.

\item \textbf{Stage 1 Gain (52$\times$):} Non-inverting MCP6002 amplifier with 1~k$\Omega$ input
  and 51~k$\Omega$ feedback resistor. Input AC-coupled via 1~\textmu F capacitor.

\item \textbf{Stage 2 Gain (21$\times$):} Inverting stage with 1~k$\Omega$ input and 21~k$\Omega$ feedback.
  Total gain: $52 \times 21 = 1092\times$ (60.8~dB).

\item \textbf{Final Anti-Aliasing Filter:} Single-pole RC with $R = 2$~k$\Omega$ and $C = 100$~nF,
  giving $f_c \approx 800$~Hz. Matches the ADC Nyquist limit (10~kHz sample rate).
\end{enumerate}

Power supply filtering uses three capacitors (0.1, 10, 100~\textmu F) with three 10~$\Omega$ series
resistors, providing broadband decoupling from DC to 10~kHz.

\subsection{Firmware and Data Acquisition}

The ESP32-S3 microcontroller (dual-core, 8~MB PSRAM) acquires 12-bit ADC samples at 10~kHz via
DMA in continuous mode. Each block contains 1024 samples from 4 channels ($\sim$102.4~ms).
The USB-CDC interface streams blocks with 0xDEADBEEF sync words and sample counts.

Protocol:
\begin{equation}
[\text{0xDEADBEEF}] [\text{count}] [\text{1024} \times 4 \times \text{uint16}]
\end{equation}

\subsubsection{Data Processing Host}

A Raspberry Pi 4 (4~GB RAM) receives blocks over USB serial at 921600~baud and runs Python
DSP pipelines using NumPy, SciPy, and custom MUSIC/Kalman filter code.

\subsection{Software Pipeline}

\begin{enumerate}
\item DC offset removal (per-channel mean subtraction)
\item Bandpass filter (4th-order Butterworth, 10--800~Hz)
\item Hanning window + 1024-point FFT
\item Peak Doppler bin detection
\item Phase extraction at peak (for MUSIC) or amplitude only (for monopulse)
\item MUSIC pseudospectrum or monopulse left/right bias computation
\item Kalman filter tracking (constant-velocity 2D motion model)
\end{enumerate}

This architecture supports two operating modes (MUSIC or monopulse) from the same data stream.

\section{Signal Processing}
\label{sec:signal}

\subsection{MUSIC Algorithm}

\subsubsection{Steering Vector and Correlation Matrix}

The MUSIC algorithm exploits the structure of the spatial correlation matrix to separate
signal and noise subspaces. For a narrowband signal at direction $\theta$ impinging on a
uniform linear array (ULA) with $M$ elements, the steering vector is:

\begin{equation}
\mathbf{a}(\theta) = [1, e^{j 2\pi d \sin(\theta) / \lambda}, e^{j 4\pi d \sin(\theta) / \lambda}, \ldots, e^{j 2\pi (M-1) d \sin(\theta) / \lambda}]^T
\label{eq:steering}
\end{equation}

where $d$ is the element spacing, $\lambda$ is the wavelength, and $\theta$ is the angle
of arrival. For our 4-element array at 10.525~GHz with $d = 37.5$~mm and $\lambda = 28.5$~mm,
the array gain is maximized at $\sin(\theta) = 0$ (broadside) and rolls off toward the edges
of the $\pm 80°$ unambiguous FOV.

The spatial correlation matrix $\mathbf{R}$ is computed over $K$ received snapshots:

\begin{equation}
\mathbf{R} = \frac{1}{K} \sum_{k=1}^{K} \mathbf{x}(k) \mathbf{x}^H(k)
\label{eq:corr_matrix}
\end{equation}

where $\mathbf{x}(k) = \mathbf{a}(\theta) s(k) + \mathbf{n}(k)$ is the received signal at time $k$,
$s(k)$ is the target signal, and $\mathbf{n}(k)$ is additive white Gaussian noise.

\subsubsection{Eigendecomposition and Pseudospectrum}

The correlation matrix is decomposed via eigendecomposition:

\begin{equation}
\mathbf{R} = \mathbf{E}_s \boldsymbol{\Lambda}_s \mathbf{E}_s^H + \mathbf{E}_n \boldsymbol{\Lambda}_n \mathbf{E}_n^H
\label{eq:eigendecomp}
\end{equation}

where $\mathbf{E}_s$ and $\mathbf{E}_n$ are the signal and noise subspace eigenvectors, respectively,
and $\boldsymbol{\Lambda}_s$, $\boldsymbol{\Lambda}_n$ are the corresponding eigenvalues. For one narrowband
signal in white noise, the noise subspace is spanned by the $M-1$ eigenvectors corresponding to the
smallest eigenvalues.

The MUSIC pseudospectrum is then:

\begin{equation}
P_{\text{MUSIC}}(\theta) = \frac{1}{\mathbf{a}^H(\theta) \mathbf{E}_n \mathbf{E}_n^H \mathbf{a}(\theta)}
\label{eq:music_pseudo}
\end{equation}

Peaks in $P_{\text{MUSIC}}(\theta)$ identify the signal directions. The algorithm achieves near-optimal
performance (close to Cramér-Rao bound) even at low SNR, provided sufficient snapshots are available~\cite{Schmidt1986}.

In practice, we search $\theta$ over the unambiguous field-of-view ($\pm 80°$) and locate the global maximum.
For multiple signals, we identify multiple peaks; however, our current hardware setup focuses on single-target
scenarios.

\subsection{Monopulse Baseline}

\subsubsection{Amplitude-Only DOA Estimation}

Monopulse is an amplitude-based direction-finding technique that divides the array into
left and right subarrays and compares the peak amplitudes:

\begin{equation}
\text{DOA\_bias} = \frac{\sum_{\text{right}} |A_{\text{right}}| - \sum_{\text{left}} |A_{\text{left}}|}{\sum_{\text{right}} |A_{\text{right}}| + \sum_{\text{left}} |A_{\text{left}}|}
\label{eq:monopulse}
\end{equation}

where $A_{\text{right}}$ and $A_{\text{left}}$ are the amplitudes (absolute values) at the peak Doppler bin
from the right and left subarrays, respectively. This bias is converted to an angle estimate via
empirical calibration or a directional response curve.

\subsubsection{Advantages and Trade-offs}

Monopulse offers three key advantages:

\begin{enumerate}
\item \textbf{Oscillator phase noise robustness:} Uses only amplitude, ignoring phase noise from free-running Gunn oscillators.
\item \textbf{Computational efficiency:} Requires only peak detection and amplitude summation (no eigendecomposition).
\item \textbf{Low latency:} Can provide DOA estimates within a single FFT frame.
\end{enumerate}

However, monopulse has fundamental limitations:

\begin{enumerate}
\item \textbf{Lower resolution:} Angular accuracy is limited by the half-power beamwidth of the subarrays,
  typically 20--30° for a 2-element subarray.
\item \textbf{Nonlinear response:} The amplitude ratio is a nonlinear function of angle; calibration is required.
\item \textbf{Sensitivity to imbalance:} Amplitude mismatches between subarrays (gain differences, insertion loss)
  directly bias the estimate.
\end{enumerate}

Despite these limitations, monopulse serves as a robust fallback when MUSIC fails under strong oscillator drift
or when real-time performance is critical.

\subsection{Extended Kalman Filter}

\subsubsection{State Model}

We employ a constant-velocity motion model in Cartesian coordinates. The state vector is:

\begin{equation}
\mathbf{x}_k = [x_k, y_k, v_{x,k}, v_{y,k}]^T
\label{eq:state_vector}
\end{equation}

where $(x_k, y_k)$ is the target position and $(v_{x,k}, v_{y,k})$ is the velocity, all at time step $k$.

The state transition model assumes constant velocity over the prediction interval $\Delta t$:

\begin{equation}
\mathbf{x}_{k+1} = \mathbf{F} \mathbf{x}_k + \mathbf{w}_k
\label{eq:state_transition}
\end{equation}

where the state transition matrix is:

\begin{equation}
\mathbf{F} = \begin{bmatrix} 1 & 0 & \Delta t & 0 \\ 0 & 1 & 0 & \Delta t \\ 0 & 0 & 1 & 0 \\ 0 & 0 & 0 & 1 \end{bmatrix}
\label{eq:transition_matrix}
\end{equation}

and $\mathbf{w}_k \sim \mathcal{N}(0, \mathbf{Q})$ is process noise representing model uncertainty.

\subsubsection{Measurement Model and Update Step}

The measurement is the estimated angle $\hat{\theta}_k$ from either MUSIC or monopulse. We convert this to
a position measurement via a range estimate (typically obtained from Doppler frequency or, if available, from
an auxiliary range sensor):

\begin{equation}
\begin{bmatrix} \hat{x}_k \\ \hat{y}_k \end{bmatrix} = r \begin{bmatrix} \sin(\hat{\theta}_k) \\ \cos(\hat{\theta}_k) \end{bmatrix} + \mathbf{v}_k
\label{eq:measurement}
\end{equation}

where $r$ is the estimated range and $\mathbf{v}_k \sim \mathcal{N}(0, \mathbf{R})$ is measurement noise.

The measurement model is nonlinear due to the trigonometric conversion; hence we apply the Extended Kalman Filter (EKF).
The measurement Jacobian with respect to the state is:

\begin{equation}
\mathbf{H} = \frac{\partial \mathbf{h}(\mathbf{x})}{\partial \mathbf{x}} = \begin{bmatrix} r \cos(\hat{\theta}_k) & 0 & \sin(\hat{\theta}_k) & 0 \\ -r \sin(\hat{\theta}_k) & 0 & 0 & \cos(\hat{\theta}_k) \end{bmatrix}
\label{eq:measurement_jacobian}
\end{equation}

At each step, the EKF performs prediction ($\mathbf{x}_{k+1|k}$) and update ($\mathbf{x}_{k+1|k+1}$) with adaptive
noise covariances. We adaptively scale $\mathbf{R}$ based on the innovation (measurement residual) to handle
oscillator-induced uncertainty.

\subsection{Forward-Backward Averaging}

\subsubsection{Motivation: Free-Running Oscillator Phase Drift}

The free-running Gunn oscillator exhibits slow phase drift over time scales of 10--100~ms. This causes the
spatial correlation matrix to become asymmetric and violates the assumption of stationarity. Forward-backward
averaging (also called conjugate symmetry averaging) exploits the real-valued nature of narrowband Doppler signals
to improve robustness.

\subsubsection{Forward-Backward Averaging Formula}

For a real-valued signal, the correlation matrix exhibits conjugate symmetry in a reference frame. Forward-backward
averaging computes:

\begin{equation}
\mathbf{R}_{\text{fb}} = \frac{1}{2} \left( \mathbf{R} + \mathbf{J} \mathbf{R}^* \mathbf{J}^T \right)
\label{eq:fb_averaging}
\end{equation}

where $\mathbf{J}$ is the exchange (reversal) matrix:

\begin{equation}
\mathbf{J} = \begin{bmatrix} 0 & 0 & 0 & 1 \\ 0 & 0 & 1 & 0 \\ 0 & 1 & 0 & 0 \\ 1 & 0 & 0 & 0 \end{bmatrix}
\label{eq:exchange_matrix}
\end{equation}

This operation enforces approximate conjugate symmetry and improves the condition number of $\mathbf{R}_{\text{fb}}$,
making the subsequent eigendecomposition more numerically stable.

\subsubsection{Benefit for Oscillator Drift}

By averaging forward and backward-reversed snapshots, we average out slow oscillator phase fluctuations that occur
between early and later samples. This reduces the effective oscillator drift over the snapshot window and improves
spatial coherence for MUSIC processing. Empirically, forward-backward averaging reduces MUSIC pseudospectrum
leakage and sharpens peaks even under 10--30~mrad/s phase drift rates~\cite{VanTrees2002}.

In our implementation, we apply forward-backward averaging to the correlation matrix before eigendecomposition,
recomputing the noise subspace with the averaged matrix. This is particularly effective for integration windows
of 100--500~ms, where oscillator drift is moderate but noticeable.

\section{Simulation \& Validation}
\label{sec:validation}

\subsection{Synthetic Signal Generation}

[Placeholder: Signal model, noise]

\subsection{Metrics and Methodology}

[Placeholder: RMS error, detection probability, Monte Carlo]

\subsection{Results}

\begin{figure}[h]
\centering
\includegraphics[width=0.8\textwidth]{figures/angle_error_vs_snr.pdf}
\caption{Angle estimation RMS error vs SNR. MUSIC outperforms monopulse across all SNR levels.}
\label{fig:angle_error}
\end{figure}

[Placeholder: Additional figures and discussion]

\section{Reproducibility}
\label{sec:repro}

[Placeholder: Code structure, Hydra configs, reproduction commands]

\section{Conclusion}
\label{sec:conclusion}

[Placeholder: Summary, future work, hardware validation]

\bibliographystyle{abbrvnat}
\bibliography{paper}

\end{document}
